\subsection{Standards Built by the Technical Community, Not by Governments}
\label{sec:8.7.2}
Section~\ref{sec:8.7.6} covers the treaty-and-declaration landscape states negotiate directly: a binding Council of Europe convention, the OECD's recommendation, UNESCO's, the two industry-and-government partnerships, the safety-institute network built after the Bletchley Declaration. A different, older layer of standard-setting sits underneath all of it, built by engineers rather than diplomats, and it deserves its own account.

The IEEE's Ethically Aligned Design, first published in 2019, set out eight general principles for autonomous and intelligent systems and, more consequentially, a follow-on series of engineering standards — the IEEE 7000 series — that translate those principles into something a development team can actually test against: value-based requirements engineering (IEEE 7000, 2021) and a well-being impact assessment (IEEE 7010), among a growing list addressing bias, transparency, and algorithmic provenance \autocite{ieee2019eadand7000series}. The document informed the OECD's own AI Principles the same year, among other governmental efforts — engineers wrote a vocabulary for what an ethical AI system does before most governments had written a rule about it \autocite{oecd2019aiprinciples}.

ISO/IEC 42001, published in December 2023, took that layer further: the first management-system standard for AI, structured like the standards a company already uses for information security or quality control, specifying what a functioning AI governance program inside an organization has to include — risk assessment, an impact assessment for each system, oversight of third-party suppliers — and auditable by an accredited certification body the way those other standards are \autocite{iso2023iec42001}.

What a technical standard does that a treaty cannot is bind through the market rather than through a state's consent. A government uninterested in constraining AI has no reason to ratify a human-rights convention; a company that wants to sell into a market, or to a customer, that requires ISO 42001 certification has a reason regardless of what its home government thinks of the standard's content. That advantage has limits — a company operating entirely inside a market that never asks for the certification has no reason to seek it either, and certification measures whether a management system exists, not whether the system it describes is followed with any rigor day to day. It is a different lever from the ones section~\ref{sec:8.7.6} covers, not a substitute for them, and it reaches an authoritarian government's own AI systems, built and used domestically with no export or customer relationship at stake, not at all.
